\IfFileExists{papers/formalized-vs-literature-preamble.tex}{\documentclass[11pt]{article}
\usepackage[utf8]{inputenc}
\usepackage[T1]{fontenc}
\usepackage{amsmath,amssymb,mathtools}
\usepackage{booktabs,tabularx,array}
\usepackage{enumitem}
\usepackage[margin=1in]{geometry}
\usepackage{xcolor}
\usepackage[hidelinks]{hyperref}
\usepackage{microtype}
\setlength{\parindent}{0pt}
\setlength{\parskip}{0.55em}
\newcommand{\Lean}[1]{\texttt{\nolinkurl{#1}}}
\newcommand{\RepoFile}[1]{\texttt{\nolinkurl{#1}}}
\newcommand{\Class}[1]{\textbf{#1}}
\newcommand{\Top}{\mathord{\top}}
\newcommand{\R}{\mathbb{R}}
\newcommand{\ENN}{\mathbb{R}_{\ge 0}^{\infty}}
\newcommand{\KL}{\operatorname{KL}}
\newcommand{\W}{\mathsf{W}}
\newcommand{\statusline}[2]{%
  \begin{center}\fbox{\begin{minipage}{0.92\linewidth}\textbf{Completion class:} #1\\[2pt]
  \textbf{Strongest caveat:} #2\end{minipage}}\end{center}}
\newenvironment{comparison}{\tabularx{\linewidth}{>{\bfseries}p{0.25\linewidth}X}}{\endtabularx}
\newcommand{\snapshot}{\texttt{902f51aa01a737af68f6feb71c83263cc9b1f4d9}}
}{\documentclass[11pt]{article}
\usepackage[utf8]{inputenc}
\usepackage[T1]{fontenc}
\usepackage{amsmath,amssymb,mathtools}
\usepackage{booktabs,tabularx,array}
\usepackage{enumitem}
\usepackage[margin=1in]{geometry}
\usepackage{xcolor}
\usepackage[hidelinks]{hyperref}
\usepackage{microtype}
\setlength{\parindent}{0pt}
\setlength{\parskip}{0.55em}
\newcommand{\Lean}[1]{\texttt{\nolinkurl{#1}}}
\newcommand{\RepoFile}[1]{\texttt{\nolinkurl{#1}}}
\newcommand{\Class}[1]{\textbf{#1}}
\newcommand{\Top}{\mathord{\top}}
\newcommand{\R}{\mathbb{R}}
\newcommand{\ENN}{\mathbb{R}_{\ge 0}^{\infty}}
\newcommand{\KL}{\operatorname{KL}}
\newcommand{\W}{\mathsf{W}}
\newcommand{\statusline}[2]{%
  \begin{center}\fbox{\begin{minipage}{0.92\linewidth}\textbf{Completion class:} #1\\[2pt]
  \textbf{Strongest caveat:} #2\end{minipage}}\end{center}}
\newenvironment{comparison}{\tabularx{\linewidth}{>{\bfseries}p{0.25\linewidth}X}}{\endtabularx}
\newcommand{\snapshot}{\texttt{902f51aa01a737af68f6feb71c83263cc9b1f4d9}}
}
\title{Finite Gaussian shift and Euler--Maruyama control-energy identity: formalized result versus the literature}
\author{}
\date{Formalization snapshot \snapshot}
\begin{document}
\maketitle
\statusline{\Class{S/V}: complete finite-dimensional model built on a vendored Gaussian-KL theorem}{This is a discrete Gaussian/Euler--Maruyama identity, not a path-space Cameron--Martin or Girsanov theorem.}
\section{Result}
The finite-dimensional Cameron--Martin entropy formula is
\[
 \KL\bigl((\mathcal N(0,I))\circ(x\mapsto x+h)^{-1}\,\Vert\,\mathcal N(0,I)\bigr)
 =\frac12\sum_i h_i^2.
\]
After scaling independent increments by $\sqrt{\Delta t}$, this yields the Euler--Maruyama identity between path-law KL and the discrete control energy $\sum_k\Delta t\,\lVert u_k\rVert^2/2$.

\section{Lean declarations}
\begin{itemize}[leftmargin=2em]
\item \Lean{ForMathlib.MeasureTheory.klDiv_stdGaussian_map_add};
\item \Lean{ForMathlib.MeasureTheory.klDiv_emShift_eq_emEnergy};
\item \Lean{ChenGeorgiouPavon2021.energy_identity_euler_maruyama}.
\end{itemize}
The first theorem is in \RepoFile{ForMathlib/MeasureTheory/GaussianEntropy.lean}; the discrete control wrapper is in \RepoFile{ChenGeorgiouPavon2021/SequenceGaussian.lean}.

\section{Provenance}
The finite Gaussian entropy block is literally vendored from Michael R. Douglas's
\url{https://github.com/mrdouglasny/gibbs-variational} at commit
\texttt{75e08d8aaca83bb090fc43855898991fbfc9abf3}, Apache-2.0, with namespace adaptation.  The Euler--Maruyama control-energy theorem is local work built on that block.

\section{Comparison}
\begin{comparison}
Dimension & Finite.\\
Covariance & Standard independent coordinates after explicit scaling.\\
Shift & Deterministic vector/control array.\\
Identity & Exact KL equals discrete energy.\\
Continuum limit & Not included.\\
Adapted random drift & Not included.\\
\end{comparison}

\section{Literature relation}
The result is the finite-coordinate atom behind Cameron--Martin and Girsanov arguments, and the local distilled sources
\RepoFile{prose/distilled_literature/CameronMartin1945_wiener_transformations.tex} and
\RepoFile{prose/distilled_literature/Girsanov1960_change_of_measure.tex} explain the continuum statements.  The Lean theorem should not be labeled as either full theorem.
\end{document}
